\chapter{Memory-Augmented 3D Object Detection}
\label{sec:method}

3D object detectors take a short temporal window of sensor data as input and produce a set of detections.
Existing approaches typically struggle to perceive occluded and distant objects with limited sensor observations.
To tackle these challenges, we propose \ourmodel{}, a plug-and-play module to enhance existing 3D object detectors with the ability to perform long-horizon temporal fusion.
Our only requirement from the detector is that each detection includes an object bounding box, multi-class confidence scores,
and a feature vector capturing local context.
We demonstrate the generality of \ourmodel{} by augmenting and improving various LiDAR-based and camera-based detectors.

We enable long-horizon temporal understanding through a memory bank that captures all the relevant information on objects,
including where we expect them to move.
These trajectory forecasts allow us to align the memory objects with the current detector proposals in space and time.
Importantly, we do not require the object detector to provide motion forecasts; instead, \ourmodel{} computes them.
To compensate for ego-motion, we assume the ego is localized --- which is the norm in modern self-driving platforms~\cite{wilson2023argoverse,WOD} --- and store the ego pose in the memory along with the model outputs.

\section{Model}
\label{sec:method-arch}
We start with an overview of our model; refer to \Cref{fig:arch} for an illustration.
At every inference step, \ourmodel{} takes as input the \textit{detection proposals}, the current timestamp $t$ (e.g., LiDAR sweep-end time or camera capture time),
and the ego pose $\mathbf{E}_{t}$ in a global coordinate frame.
It then retrieves objects from memory, aligns them spatially with $\mathbf{E}_t$ and temporally to $t$, and extracts high-dimensional features from the aligned boxes and trajectory forecasts.
We refer to the aligned boxes and trajectory forecasts
with the extracted features as \textit{memory proposals}.
A proposal merging mechanism then fuses detection and memory proposals by rescoring their confidence scores and applying standard post-processing.
Finally, our refinement transformer iterative refines the object detections and trajectory forecasts in the \textit{merged proposals} with cross-attention to the memory and factorized self-attention.
In preparation for future inferences, the memory bank is then updated by appending the model outputs (a.k.a. \textit{refined proposals}) and removing older model outputs to keep the memory bounded in size. 

\subsection{Proposal representation}  
We define object \textit{proposals}
$\mathcal{P} = (\mathbf{B}, \mathbf{C}, \mathbf{T}, \mathbf{Q})$ with
$N$ bounding boxes $\mathbf{B} \in \mathbb{R}^{N \times 7}$,
where the last dimension corresponds to $(x, y, z, l, w, h, \theta)$ with
object 3D centroids $(x, y, z)$,
headings $\theta$ in a BEV ego-relative coordinate frame,
and the 3D box dimensions $(w, l, h)$;
multi-class confidence scores $\mathbf{C} \in [0, 1]^{N \times C}$,
where $C$ is the number of actor classes;
trajectory forecasts $\mathbf{T} \in \mathbb{R}^{N \times T_f \times 3}$ describing objects' BEV pose
$\{(x, y, \theta)_{t+s_f}, \dots, (x, y, \theta)_{t + T_f s_f}\}$
over $T_f$ future waypoints at a time interval $s_f$;
and an object feature $\mathbf{Q} \in \mathbb{R}^{N \times (T_f + 1) \times d}$ encoding both local and global features for every object at the present and future timestamps,
where $d$ is the feature dimensionality.
We use superscripts to denote the source of the proposals:
detection proposals $\mathcal{P}^\mathrm{det}$ from the 3D detector,
memory proposals $\mathcal{P}^\mathrm{mem}$ from the memory bank,
merged proposals $\mathcal{P}^\mathrm{merge}$ from the proposal merging module,
and refined proposals $\mathcal{P}^\mathrm{ref}$ from the output of the refinement transformer.
For detection proposals $\mathcal{P}^{\mathrm{det}}$,
we generate $\mathbf{T}^\mathrm{det}$
by assuming the object is static over time since detectors do not provide forecasts (and this is just an initialization before refinement). The object features $\mathbf{Q}^\mathrm{det}$ are obtained by interpolating the feature map from before the detector header at the projected object centroids,
repeating $(T_f + 1)$ times to get the features for future timestamps, and adding a learned embedding of $\mathbf{B}^\mathrm{det}$ and $\mathbf{C}^\mathrm{det}$.
In the paragraphs below we describe how we obtain $\mathcal{P}^\mathrm{mem}$, $\mathcal{P}^\mathrm{merge}$ and $\mathcal{P}^\mathrm{ref}$.

\subsection{Memory Bank and Retrieval}
The memory bank is a set of tuples $\big(t_m, \mathbf{E}_{t_m}, \mathcal{P}^{\mathrm{ref}}_{t_m} \big)$
with timestamped past model outputs and ego pose, sorted by the timestamp $t_m$ at which the outputs were generated.
During inference at timestamp $t$, we retrieve memory entries
$\mathcal{P}_{t_m}^\mathrm{ref}$
at a set of past target timestamps $t_m \in \mathcal{T}_m$, where $\mathcal{T}_m = \{t- s_m, t-2 s_m, \dots, t-T_m s_m\}$.
$T_m$ is the number of past target timestamps, and $s_m$ is the time stride of the retrieved entries. To be precise, we retrieve the closest memory entry to each timestamp in $\mathcal{T}_m$ to be robust to small sensor delays.

\subsection{Extracting Memory Proposals}
For effective use of the memory at inference,
we should align each retrieved entry
$(t_m, \mathbf{E}_{t_m}, \mathcal{P}^{\mathrm{ref}}_{t_m})$
in space and time with the current detection proposals at time $t$.
We handle ego-motion by applying the relative transform
$\mathbf{E}_{t_m \to t} = \mathbf{E}_{t}^{-1} \mathbf{E}_{t_m}$
to $\mathbf{B}_{t_m}^{\mathrm{ref}}$ and $\mathbf{T}_{t_m}^{\mathrm{ref}}$.
To handle object motion, we linearly interpolate the stored trajectory forecast to the current timestamp $t$ to obtain the proposal box $\mathbf{B}^\mathrm{mem}$.
To obtain the proposal forecast $\mathbf{T}^\mathrm{mem}$,
we also interpolate/extrapolate the stored trajectories as required to obtain waypoints at $\mathcal{T}_f = \{t+s_f, \dots, t + T_f s_f\}$ from stored waypoints at $\{t_m+s_f, \dots, t_m + T_f s_f\}$.

Finally, we extract latent features $\mathbf{Q}^\mathrm{mem}$ at
$t$ and every future time step $t_f \in \mathcal{T}_f$:
First, we compute sinusoidal positional embeddings~\cite{vaswani2017attention}
for the centroid coordinates $\mathbf{B}^{\mathrm{mem}}_{x, y, z}$ and encode
them with a lightweight MLP.
Separately, we concatenate other features including
$\mathbf{B}^{\mathrm{mem}}_{l, w, h, \theta}$ (box dimensions and heading), confidence scores $\mathbf{C}^{\mathrm{mem}}$, the memory age $t - t_m$,
and a 2D vector pointing to where the proposal was in the current ego coordinate frame at the time $t_m$.
Finally, we encode the concatenated features with another MLP and add the features from both MLPs together.

\subsection{Proposal Merging} 
The memory and detection proposals can be redundant, particularly in areas with good sensor coverage. 
To merge proposals, we learn to rescore their multi-class confidence scores.
Rescoring is essential as the confidence the model should put in a memory proposal
not only depends on the confidence score at a past timestamp
$\mathbf{C}_{t_m}^\mathrm{mem}$, but also on the proposal age $t - t_m$, as the forecasting uncertainty grows with the time horizon and other factors.
For example, the model should trust a fast-moving detection less than a stationary object, or it should trust an object observed 0.5 seconds ago more than one observed 5 seconds ago.
Furthermore, the detection proposals come from the 3D detector, while the memory proposals are produced by \ourmodel{},
and detectors have been found to be miscalibrated~\cite{neumann2018relaxed,kuppers2020multivariate,minderer2021revisiting}.

To make the scores comparable, we learn two small MLPs that separately map the features of the detection proposals $\mathbf{Q}^\mathrm{det}$ and the memory proposals $\mathbf{Q}^\mathrm{mem}$
to new multi-class scores $\mathbf{C}^\mathrm{merge}$.
As explained in \Cref{sec:method-training},
these rescoring MLPs are trained under a single detection loss applied to the merged proposals so that the model can decide which proposals to trust from both sources.
Finally, we filter proposals with score thresholding, non-maximum suppression (NMS), and keep the top $K$
\textit{merged proposals} sorted by score (maximum over actor classes).
Post-processing enables the refinement transformer to process a smaller number of queries.

Finally, we add learned time positional embeddings to the merged proposal features $\mathbf{Q}^\mathrm{merge}$
to indicate the time of the trajectory forecast.
At this point, we have $N^\mathrm{merge} \overset{\mathrm{def}}{=} N^\mathrm{ref}$ merged proposals $\mathcal{P}^\mathrm{merge}$ ready for refinement.

\subsection{Refinement Transformer}
We utilize a transformer decoder to refine the merged proposals $\mathcal{P}^\mathrm{merge} = \mathcal{P}^{\mathrm{ref}(0)}$
iteratively over $I$ blocks into $\mathcal{P}^{\mathrm{ref}(1)}$ \dots $\mathcal{P}^{\mathrm{ref}(I)}$,
where the final model outputs are $\mathcal{P}^\mathrm{ref} = \mathcal{P}^{\mathrm{ref}(I)}$.
We propose a novel \textit{memory cross-attention}
mechanism to allow the queries --- proposal features $\mathbf{Q}^{\mathrm{ref}(i)}$ --- to aggregate information from all the memory proposals $\mathbf{Q}^{\mathrm{mem}}$,
including those that proposal merging filtered out.
We want to use this information in the refinement transformer because multiple overlapping memory proposals provide significant evidence about an object's presence and location.
To achieve this, we perform cross attention from the object queries $\mathbf{Q}^{\mathrm{ref}(i)}$ to the memory proposal features $\mathbf{Q}^\mathrm{mem}$.
For efficiency, we limit the cross attention to the nearest $k$ keys to each object query (computing the nearest neighbors of $\mathbf{B}^{\mathrm{ref}(i)}_{x, y, z}$ in $\mathbf{B}^{\mathrm{mem}}_{x, y, z}$).

Similar to many works~\cite{ngiam2021scene,casas2024detra}, we perform factorized self-attention in each refinement block, which separates \textit{time self-attention} and \textit{object self-attention} for efficiency, where the former attends only to queries from the same object (sequence length $T_f + 1$) and the latter only attends to queries from the same time step (sequence length $N$).
The updated queries $\mathbf{Q}^{\mathrm{ref}(i+1)}$ are input to the next block.

Finally, we update the explicit proposal information as described in DeTra~\cite{casas2024detra}, by using a simple MLP to produce $\mathbf{B}^{\mathrm{ref}(i+1)}$ and $\mathbf{C}^{\mathrm{ref}(i+1)}$ and a gated recurrent unit (GRU)~\cite{cho2014learning} to update the future trajectory waypoints $\mathbf{T}^{\mathrm{ref}(i+1)}$.

\subsection{Memory Bank Update}
We post-process the refined proposals $\mathcal{P}^{\mathrm{ref}}$ as we did to the merged proposals: score thresholding,
NMS, and top $K$ based on confidence score,
adding the result to the memory bank, along with the corresponding timestamp $t$ and ego pose $\mathbf{E}_t$.
To limit the size of the memory bank when running on long sequences,
we remove any memory entries older than $t - T_m s_m - \epsilon$ (the past time-horizon used in memory retrieval with a small buffer $\epsilon$).